% ============================================================
% sec/90respaldo.tex — diapositivas de respaldo (no cuentan)
% El mapa IN debe coincidir con el PDF del atributo IN.
% ============================================================

\begin{frame}{Respaldo: mapa del atributo IN}
  \decidir{alinear cada fila con el texto literal de IN1--IN5}

  \vspace{0.1cm}
  \scriptsize
  \begin{tabularx}{\textwidth}{@{}l P{2.6cm} Y P{3.0cm}@{}}
    \toprule
    \textbf{Indicador} & \textbf{Tema} & \textbf{Diapositivas} &
    \textbf{Informe}\\
    \midrule
    IN1 & Problema & Contexto: pipeline y tr\'iada CID & Secci\'on I\\
    IN2 & B\'usqueda & Fuentes, palabras clave y embudo de selecci\'on &
      Secci\'on II\\
    IN3 & An\'alisis & S\'intesis cruzada: tres ejes y contradicciones &
      Secciones III--IV, Tabla I\\
    IN4 & Evidencia & Demostraci\'on: montaje, umbral y resultados &
      Carpeta \texttt{demo/}\\
    IN5 & Conclusiones & Conclusiones y qu\'e vigilar & Secciones V--VI\\
    \bottomrule
  \end{tabularx}
\end{frame}

\begin{frame}[allowframebreaks]{Referencias}
  \tiny
  \bibliographystyle{ieeetr}
  \bibliography{sec/99referencias}
\end{frame}